% Unification: knowability = margin > reach. Numbers from
% scripts/unification_reach_validation.py (results/theory/unification_reach_validation.json).

The per-family boundaries proved above share one shape. This subsection makes that a
theorem: a single computable quantity --- the \emph{ambiguity reach} --- decides
knowability in every family in this paper.

\begin{definition}[Evidence class, benefit interval, reach]\label{def:reach}
Let $\mathcal F$ be a family of admissible worlds and $O(P)$ the joint law of all
label-free observables of world $P$ (unlabeled target data and the labeled source).
The \emph{evidence class} of $P$ is $[P]=\{Q\in\mathcal F: O(Q)=O(P)\}$, its
\emph{benefit interval} $I(P)=\{\Delta(Q):Q\in[P]\}$, with \emph{center}
$c(P)=\tfrac12(\sup I+\inf I)$ and \emph{reach} $\rho(P)=\tfrac12(\sup I-\inf I)$.
\end{definition}

\begin{proposition}[Knowability $=$ margin beyond reach]\label{thm:unify}
Assume the regularity needed for a measurable selection of a point of $[P]$
($\mathcal F$ Polish, $O$ and $\Delta$ measurable, $[P]$ compact-valued; standard and
flagged as an assumption). Then:
\textup{(i)} if $|c(P)|>\rho(P)$ for all $P\in\mathcal F$, the rule that estimates the
center by any consistent $\widehat c$ with radius $\varepsilon_n$ and commits
$\sign\widehat c$ iff $|\widehat c|>\rho+\varepsilon_n$ is $\alpha$-valid over the whole
family;
\textup{(ii)} conversely, if $0$ lies in the interior of $I(P)$ for some $P$, two worlds
of $[P]$ carry opposite signs of $\Delta$ with identical label-free observables ---
Theorem~\ref{thm:imp} applies and no label-free rule with wrong-commit
$<\tfrac12$ can commit on $[P]$.
Hence $\sign\Delta$ is identifiable on $\mathcal F$ exactly where the margin $|c|$
exceeds the reach $\rho$.
\end{proposition}
\begin{proof}
(i) On the $1-\alpha$ event, $|\widehat c-c|\le\varepsilon_n$, so
$|\widehat c|>\rho+\varepsilon_n$ implies $|c|>\rho$, and every $Q\in[P]$ has
$\Delta(Q)\in[c-\rho,c+\rho]$, an interval not containing $0$, of the committed sign.
(ii) Both signs are attained on $I(P)$ by definition of interior; pick two attaining
worlds; their observable laws coincide by membership in $[P]$, and
Theorem~\ref{thm:imp}'s argument (identical action law across the pair) forbids a valid
commitment.
\end{proof}

\begin{proposition}[The reach table]\label{prop:reach-table}
The boundaries of this paper are evaluations of $\rho$:
\begin{center}
%% fix-queue item 32: \resizebox dropped; set at \small instead of rescaled.
\small
{%
\begin{tabular}{llll}
\hline
family & center $c$ & reach $\rho$ & source \\
\hline
covariate shift & $\Delta$ & $0$ & Thm.~\ref{thm:pos} \\
bounded drift $\|b\|_\infty\!\le\!B$ & $U-2T_S$ & $2B\,\E_T|f_0{-}f_a|$ & Prop.~\ref{thm:reg-iff} \\
label shift & midpoint over $A(\pi)$ & $\sup_{A(\pi)}|\delta\cdot v|$ & Prop.~\ref{thm:ls-iff} \\
multi-view (cond.\ indep.) & $\Delta$ & $0$ under moment non-degeneracy & known (Steinhardt--Liang) \\
two-point witness & $0$ & $|\Delta|$ & Thm.~\ref{thm:imp} \\
\hline
\end{tabular}}
\end{center}
Each row follows by computing $I(P)$ in the family; the witness row shows the maximal
case $c=0$, $\rho=|\Delta|$ --- never identifiable. The margin $\kappa_\alpha$ of
Proposition~\ref{prop:kfrontier} is the finite-sample shadow of $|c|-\rho$.
\end{proposition}

\noindent\emph{Numerically validated}
(\texttt{unification\_reach\_validation.py}; Fig.~\ref{fig:reachtable}). Brute-force
suprema over admissible equivalent worlds reproduce the closed-form reaches with ratio
$1.00$ in both nontrivial families (bounded drift; singular label shift); the
margin-beyond-reach certificate makes zero false certifications in all three families;
and sign flips among equivalent worlds occur exactly when the margin is below the
reach.

\begin{figure}[t]\centering
\includegraphics[width=0.8\linewidth]{figures/fig_reach_table.png}
\caption{One quantity decides every family (Proposition~\ref{thm:unify}): identifiable iff
the margin $|c|$ exceeds the reach $\rho$. Bars: closed forms; diamonds: brute-force
empirical reach over admissible equivalent worlds.}
\label{fig:reachtable}
\end{figure}

\paragraph{Weight (honest).} Direction (ii) is a corollary of
Theorem~\ref{thm:imp}; direction (i) is close to definitional once the stated
selection regularity is assumed --- we say so plainly. The contribution is the
\emph{organization}: a single computable quantity whose per-family evaluations are
exactly this paper's positive and negative results, and a sharp restatement of the open
problem: the general case of Conjecture~\ref{conj:gen} \emph{is} the computation of
$\rho$ for nonparametric families (unknown conditionals, unbounded drift), where no
closed form is known.
